%! Introducing Statistics: What Can We Learn from Data? — Unit 1, Lesson 1.0 — Guided Notes & Practice
% GRADUAL RELEASE component (I do -> we do -> you do), 22 minutes.
%   I DO   : the objective box, the vocabulary box, and notes sections 1-4. The teacher
%            presents; students fill the blanks as each idea is named.
%   WE DO  : the practicebox ("Guided Practice") — worked together, students hold the pen.
%   YOU DO : the "Independent Practice" notesbox — students work alone, teacher circulates.
% Vocabulary is GIVEN here, not discovered: the group activity that follows applies it.
%
% PAGE LOCKSTEP with notes_key/: every \blank{} in this file becomes a short \ans{} there, and
% every \writelines{n} becomes exactly n \ansline{} calls. Keep key answers short enough that
% they do not wrap past the reserved lines, and keep every \boxguard byte-identical.
% VOCABPAR: \par brackets each \termblank inside the vocabbox — without it the intro sentence
% collides with the first term and key answers drag onto the wrong line.
\documentclass[10pt]{article}
\usepackage{apstats-article}
\usepackage{apstats-boxes}

\begin{document}

\pageheader{Unit 1, Lesson 1.0}{Guided Notes \& Practice}
% NAMESTRIP: no \namedateperiod here — the name row lives on the cover only.

\vspace{0.12in}

\boxguard[14]
\begin{objectivebox}
\small
By the end of these notes I will be able to:
\begin{enumerate}[leftmargin=1.6em, nosep, after=\vspace{2pt}]
  \item explain why a number carries no information until it is placed in \textbf{context};
  \item identify the \textbf{individuals}, \textbf{variables}, and \textbf{values} in a data set;
  \item decide whether a question is a \textbf{statistical question} and justify it by naming
        what varies.
\end{enumerate}
\end{objectivebox}

\vspace{0.10in}

% ── VOCABULARY (I do — students fill each term as it is named) ────────────────
\boxguard[20]
\begin{vocabbox}
\small
Fill in each term \textbf{as we name it} in the notes below.
\par
\termblank{Individual}
\par
\termblank{Variable}
\par
\termblank{Value}
\par
\termblank{Variability}
\par
\termblank{Statistical question}
\par
\end{vocabbox}

\vspace{0.10in}

% ── I DO: section 1 ──────────────────────────────────────────────────────────
\boxguard[16]
\begin{notesbox}{1. A Number Is Not Data Until It Has Context}
\small
The seven numbers from the warm-up --- 68, 71, 64, 70, 66, 73, 69 --- told you nothing. They
could have been test scores, temperatures, ages, or heights. Nothing changes about the numbers
themselves; what was missing was the \textbf{context}.

\vspace{6pt}
\textbf{Statistics} is the science of \blank{2.4cm}, organizing, analyzing, and drawing
\blank{2.6cm} from data.

\vspace{4pt}
\textbf{Data} are recorded values \emph{together with} the \blank{2.4cm} that gives them
meaning. \emph{Notice what is not in either definition: arithmetic.}

\vspace{8pt}
\textbf{The W questions.} Answer these and a bare number becomes data.

\vspace{4pt}
\renewcommand{\arraystretch}{1.45}
\begin{tabularx}{\linewidth}{@{} >{\bfseries\color{navy}}l @{\hspace{0.8cm}} X @{}}
\hline
Who?           & \blank{9.5cm} \\
What?          & \blank{9.5cm} \\
In what units? & \blank{9.5cm} \\
When \& where? & \blank{9.5cm} \\
\hline
\end{tabularx}

\vspace{6pt}
\textbf{Essential Knowledge VAR-1.A.1:} numbers convey meaningful information only when placed
in context. This is the standard every answer you write this year is held to.
\end{notesbox}

\vspace{0.10in}

% ── I DO: section 2 ──────────────────────────────────────────────────────────
\boxguard[24]
\begin{notesbox}{2. Individuals, Variables, and Values}
\small
A survey of four students in a first-period class:

\vspace{5pt}
\renewcommand{\arraystretch}{1.25}
\begin{tabularx}{\linewidth}{@{} l Y Y Y @{}}
\rowcolor{sky}
\textbf{Student} & \textbf{Grade level} & \textbf{Hours of sleep} & \textbf{Rides the bus?} \\
\hline
Alvarez & 10 & 7.5 & Yes \\
Boone   & 11 & 6.0 & No  \\
Chen    & 10 & 8.5 & Yes \\
Delgado & 12 & 5.5 & No  \\
\hline
\end{tabularx}

\vspace{8pt}
\begin{itemize}[leftmargin=1.3em, itemsep=6pt]
  \item An \textbf{individual} is what one \blank{2.0cm} of the table describes. Here the
        individuals are \blank{4.4cm}.
  \item A \textbf{variable} is a characteristic recorded for every individual --- a
        \blank{3.0cm} heading. This table records \blank{1.0cm} variables.
  \item A \textbf{value} is one \blank{2.0cm} underneath a variable. ``8.5'' is a value
        of the variable \blank{3.2cm}.
\end{itemize}

\vspace{6pt}
\begin{center}
\fcolorbox{redacc}{redbg}{\parbox{0.94\linewidth}{\small
\textbf{The day-one error:} naming a \emph{value} where a \emph{variable} belongs. ``Yes'' is
not a variable --- \emph{Rides the bus?} is. Before you answer, point at the column
\blank{2.6cm}.}}
\end{center}
\end{notesbox}

\vspace{0.10in}

% ── I DO: section 3 ──────────────────────────────────────────────────────────
\boxguard[16]
\begin{notesbox}{3. Variability --- Why This Course Exists}
\small
Cover the \emph{Hours of sleep} column with your hand and imagine every entry read 7.0.

\vspace{6pt}
\textbf{Variability} means the values of a variable \blank{2.6cm} from one individual to the
next.

\vspace{6pt}
If every individual gave the \blank{2.4cm} value, one measurement would answer everything ---
no sample, no graph, no average, no course. Because the values differ (7.5, 6.0, 8.5, 5.5), one
measurement is \blank{2.6cm}.

\vspace{4pt}
\textbf{Variability is the reason statistics exists.} Every technique in this course is a way
of describing, measuring, or accounting for it.
\end{notesbox}

\vspace{0.10in}

% ── I DO: section 4 ──────────────────────────────────────────────────────────
\boxguard[24]
\begin{notesbox}{4. What Makes a Question Statistical?}
\small
A \textbf{statistical question} anticipates variability and is answered by collecting data. It
must pass \textbf{both} checks:

\vspace{6pt}
\begin{tabularx}{\linewidth}{@{} >{\bfseries\color{navy}}l X @{}}
Check 1 & Do the answers \blank{2.4cm} across individuals? \\[4pt]
Check 2 & Must I \blank{2.6cm} data to answer it? \\
\end{tabularx}

\vspace{8pt}
\textbf{Careful:} needing the whole table is \emph{not} the test, and having a number in the
answer is \emph{not} the test. The test is \textbf{variation across individuals}.

\vspace{6pt}
Write \textbf{S} or \textbf{N} for each, then name what varies (or why nothing does).

\vspace{4pt}
\renewcommand{\arraystretch}{1.5}
\begin{tabularx}{\linewidth}{@{} X c X @{}}
\hline
\textbf{Question} & \textbf{S/N} & \textbf{What varies?} \\
\hline
How many students are in this table?          & \blank{0.8cm} & \blank{4.2cm} \\
How many hours did Chen sleep?                & \blank{0.8cm} & \blank{4.2cm} \\
How do the hours of sleep vary in this class? & \blank{0.8cm} & \blank{4.2cm} \\
What grade levels are represented here?       & \blank{0.8cm} & \blank{4.2cm} \\
\hline
\end{tabularx}

\vspace{6pt}
\emph{Look at the last row.} Its answer is not a number --- and it is still statistical, because
grade level differs from student to student.
\end{notesbox}

\vspace{0.10in}

% ── WE DO: guided practice ───────────────────────────────────────────────────
\boxguard[20]
\begin{practicebox}
\small
\textbf{The Gym Sign-Up --- we work this one together.} A gym records the \textbf{resting heart
rate}, in beats per minute, of each of its 120 members on the morning they join. It also records
each member's \textbf{home ZIP code}.

\begin{enumerate}[label=\textbf{\alph*.}, leftmargin=1.8em, itemsep=7pt]
  \item The individuals are \blank{4.6cm}, and the two variables recorded are
        \blank{3.4cm} and \blank{3.0cm}.
  \item Write one statistical question this gym could answer, and state exactly what varies
        in it.\\[3pt]
        \writelines{2}
  \item ``What is the most common ZIP code among the members?'' --- statistical or not? Run
        \textbf{both} checks before you answer.\\[3pt]
        \writelines{2}
\end{enumerate}
\end{practicebox}

\vspace{0.10in}

% ── YOU DO: independent practice ─────────────────────────────────────────────
\boxguard[20]
\begin{notesbox}{Independent Practice --- You Try}
\small
Work these on your own. I will circulate.

\begin{enumerate}[leftmargin=1.9em, itemsep=8pt]
  \item A biologist weighs each of the 30 turtles in a pond and records the weight in kilograms.

        The individuals are \blank{4.4cm}; the variable is \blank{3.0cm},
        measured in \blank{2.2cm}.
  \item \textbf{Same pond, two questions.} Exactly one is statistical. Circle it, then explain
        which check the other one fails.

        \vspace{4pt}
        \hspace{1.0em}\emph{(i)} How much does the heaviest turtle weigh?
        \hfill
        \emph{(ii)} How do the turtles' weights vary?
        \hspace{1.0em}\\[5pt]
        \writelines{2}
  \item A report states only: \textbf{``The average is 12.''} Name \textbf{two} pieces of
        context you would need before this number means anything.\\[3pt]
        \writelines{2}
\end{enumerate}
\end{notesbox}

\end{document}
